% ============================================================
% Section 4: Experimental Results
% ============================================================
\section{Experimental Results}

\begin{frame}{Experiment Setup}
    \textbf{Dataset:} Synthetic 2D Blobs with varying characteristics

    \vspace{0.3cm}

    \textbf{Three Experimental Scenarios:}
    \begin{enumerate}
        \item \textbf{Varying Data Distributions:} Different numbers of points and cluster densities
        \item \textbf{Kernel Function Comparison:} Gaussian, Flat (Uniform), Epanechnikov
        \item \textbf{Bandwidth Selection:} Small $h$, Optimal $h$, Large $h$
    \end{enumerate}

    \vspace{0.3cm}

    \textbf{Evaluation Metrics:}
    \begin{itemize}
        \item Silhouette Score (cluster quality)
        \item Adjusted Rand Index (vs. ground truth)
        \item Number of clusters found
        \item Computation time
    \end{itemize}
\end{frame}

\begin{frame}{Experiment 1: Varying Data Distributions}
    \textbf{Setup:} Test Mean Shift on different blob configurations

    \vspace{0.3cm}

    \begin{center}
        \textit{[Insert visualization: 3 different blob configurations]}

        \vspace{0.2cm}

        \small{(a) Sparse (100 pts) | (b) Medium (500 pts) | (c) Dense overlapping (1000 pts)}
    \end{center}

    \vspace{0.3cm}

    \begin{table}
        \centering
        \small
        \begin{tabular}{lcccc}
            \toprule
            \textbf{Configuration} & \textbf{Clusters Found} & \textbf{Silhouette} & \textbf{ARI} & \textbf{Time (s)} \\
            \midrule
            Sparse (100 pts) & 3 & 0.68 & 0.85 & 0.08 \\
            Medium (500 pts) & 3 & 0.72 & 0.91 & 0.45 \\
            Dense (1000 pts) & 4 & 0.65 & 0.78 & 1.23 \\
            \bottomrule
        \end{tabular}
    \end{table}

    \vspace{0.2cm}

    \small{\textbf{Observation:} Performance degrades with overlapping clusters and increases computation time with more points}
\end{frame}

\begin{frame}{Experiment 2: Kernel Function Comparison}
    \textbf{Setup:} Same dataset (500 points, 3 clusters), different kernels

    \vspace{0.3cm}

    \begin{center}
        \textit{[Insert visualization: Side-by-side kernel comparison]}

        \vspace{0.2cm}

        \small{Gaussian | Flat (Uniform) | Epanechnikov}
    \end{center}

    \vspace{0.3cm}

    \begin{table}
        \centering
        \begin{tabular}{lcccc}
            \toprule
            \textbf{Kernel} & \textbf{Clusters Found} & \textbf{Silhouette} & \textbf{ARI} & \textbf{Time (s)} \\
            \midrule
            Gaussian & 3 & 0.74 & 0.92 & 0.58 \\
            \textcolor{customred}{Flat (Uniform)} & \textcolor{customred}{3} & \textcolor{customred}{0.72} & \textcolor{customred}{0.91} & \textcolor{customred}{0.31} \\
            Epanechnikov & 3 & 0.73 & 0.91 & 0.42 \\
            \bottomrule
        \end{tabular}
    \end{table}

    \vspace{0.2cm}

    \small{\textbf{Key Finding:} Flat kernel offers best speed-accuracy trade-off; results generally similar across kernels}
\end{frame}

\begin{frame}{Experiment 3: Bandwidth Sensitivity}
    \textbf{Setup:} Same dataset, varying bandwidth parameter $h$

    \vspace{0.3cm}

    \begin{center}
        \textit{[Insert visualization: Clustering with different bandwidths]}

        \vspace{0.2cm}

        \small{Small $h = 0.3$ | Optimal $h = 0.8$ | Large $h = 2.0$}
    \end{center}

    \vspace{0.3cm}

    \begin{table}
        \centering
        \begin{tabular}{lcccc}
            \toprule
            \textbf{Bandwidth $h$} & \textbf{Clusters Found} & \textbf{Silhouette} & \textbf{ARI} & \textbf{Iterations} \\
            \midrule
            Small ($h = 0.3$) & 12 & 0.41 & 0.23 & 8.2 \\
            \textcolor{customred}{Optimal ($h = 0.8$)} & \textcolor{customred}{3} & \textcolor{customred}{0.72} & \textcolor{customred}{0.91} & \textcolor{customred}{4.5} \\
            Large ($h = 2.0$) & 1 & 0.18 & 0.00 & 2.1 \\
            \bottomrule
        \end{tabular}
    \end{table}

    \vspace{0.2cm}

    \small{\textbf{Critical Insight:} Bandwidth is the most important parameter}
    \begin{itemize}
        \small
        \item Too small → over-segmentation (12 clusters instead of 3)
        \item Too large → all points merge into 1 cluster
    \end{itemize}
\end{frame}

\begin{frame}{Bandwidth Selection Strategies}
    \textbf{Challenge:} How to choose optimal bandwidth $h$?

    \vspace{0.5cm}

    \begin{columns}[T]
        \begin{column}{0.5\textwidth}
            \textbf{Common Approaches:}
            \begin{enumerate}
                \item \textbf{Scott's Rule:}
                \[
                h \approx n^{-1/(d+4)} \cdot \sigma
                \]
                where $\sigma$ is std. dev.

                \vspace{0.3cm}

                \item \textbf{Silverman's Rule:}
                \[
                h \approx 1.06 \cdot \sigma \cdot n^{-1/5}
                \]

                \vspace{0.3cm}

                \item \textbf{Cross-Validation:}
                Grid search with validation metric
            \end{enumerate}
        \end{column}

        \begin{column}{0.5\textwidth}
            \begin{center}
                \textit{[Insert plot: Number of clusters vs. $h$]}

                \vspace{0.3cm}

                \small{Bandwidth $h$ vs. Clusters found}
            \end{center}

            \vspace{0.5cm}

            \textbf{Practical Tip:}
            \begin{itemize}
                \item Start with Scott's/Silverman's rule
                \item Fine-tune based on domain knowledge
                \item Use silhouette score to validate
            \end{itemize}
        \end{column}
    \end{columns}
\end{frame}

\begin{frame}{Summary of Experimental Findings}
    \textbf{Key Takeaways:}

    \vspace{0.3cm}

    \begin{enumerate}
        \item \textbf{Data Characteristics:}
        \begin{itemize}
            \item Works well on well-separated clusters
            \item Struggles with heavy overlap
            \item Computation time scales with data size
        \end{itemize}

        \vspace{0.3cm}

        \item \textbf{Kernel Choice:}
        \begin{itemize}
            \item Flat/Uniform kernel: fastest, good results
            \item Gaussian: slightly better quality, slower
            \item Choice matters less than bandwidth selection
        \end{itemize}

        \vspace{0.3cm}

        \item \textbf{Bandwidth Selection:}
        \begin{itemize}
            \item \textcolor{customred}{Most critical parameter}
            \item Use adaptive methods (Scott's/Silverman's rule)
            \item Validate with silhouette score or domain knowledge
        \end{itemize}
    \end{enumerate}
\end{frame}
